\appendix

\section{The first and second variation formulas}

We show the first and second variation formulae for finite weighted graphs into Riemannian manifolds. This result is proved by Kotani and Sunada~\cite[Theorem~2.1]{KotaniSunadaStandard}; we give a proof for the sake of convenience.


Let $X=(V, E, m_E)$ be a finite weighted graph and $(M, G)$ a Riemannian manifold. A map $f: X\to M$ is called \emph{piecewise $C^k$} if $f_{e}:=f|_e: [0,1]\to M$ is $C^k$ for each edge $e\in E$. We assume $k\geq 2$ throughout this appendix. Consider a smooth family of piecewise $C^k$-maps $f_{s}: X \to (M, G)$ for $s\in (-\epsilon, \epsilon)$ with $f:=f_{0}$ for $\e>0$. Let $f_e(s, t):=f_{s,\,e} (t)$ and $T_e(s, t):=(\p/\p t)f_e(s, t)$ for each edge $e\in E$. We denote the variation vector fields by $(\p/\p s)f_e(s, t)=V_e(s, t)$. We denote the Levi--Civita connection and its curvature tensor of $G$ by $\nabla$ and $R$, respectively.


\begin{lemm}\label{Lem:first and second}
For any $s \in (-\e, \e)$, we have the first variation formula:
\begin{multline}\label{Eq:firstVF}
\frac{d}{d s} E_G(f_{s})\\
=-2\sum_{e\,\in\,E}m_E(e)G(V_e, T_e)(s,0)-\sum_{e\,\in\,E}m_E(e)\int_0^1 G\left(V_e, \nabla_{T_e}T_e\right)(s,t) dt,
\end{multline}
and the second variation formula:
\begin{multline}\label{Eq:secondVF}
\frac{d^2}{d s^2} E_G(f_{s})=-2\sum_{e\,\in\,E}m_E(e)G\left(\nabla_{V_e} V_e, T_e\right)(s,0)\\
+\sum_{e\,\in\,E}m_E(e) \int_0^1 \big\{-G\left(\nabla_{V_e}V_e, \nabla_{T_e}T_e\right)+\left\|\nabla_{T_e}V_e\right\|_G^2 -G\left(R\left(V_e, T_e\right)T_e, V_e\right) \big\}(s, t)dt.
\end{multline}
\end{lemm}

\begin{proof}
First, we see that
\begin{equation}\label{f1}
\begin{aligned}
\frac{d}{d s} E_G(f_{s}) &=\frac{1}{2}\sum_{e\,\in\,E}m_E(e)\int_0^1 \frac{\p}{\p s}G\left(T_e, T_e\right)\,dt =\sum_{e\,\in\,E}m_E(e)\int_0^1 G\left(\nabla_{V_e}T_e, T_e\right)\, dtr\\
&=\sum_{e\,\in\,E}m_E(e)\int_0^1G\left(\nabla_{T_e}V_e, T_e\right)\, dt \qquad  \left(\text{use }[V_e, T_e]=0\right)\\
&=\sum_{e\,\in\,E}m_E(e)\int_0^1\frac{\p}{\p t}G(V_e, T_e)- G\left(V_e, \nabla_{T_e}T_e\right)\, dt\\
&=\sum_{e\,\in\,E}m_E(e)\big[G\left(V_e, T_e\right)\big]_{t=0}^{t=1}-\sum_{e\,\in\,E}m_E(e)\int_0^1 G\left(V_e, \nabla_{T_e}T_e\right)\, dt.
\end{aligned}
\end{equation}
Since
\begin{equation}\label{Eq:ap1}
V_e(s, 1)=V_{\wbar e}(s, 0) \quad \text{and} \quad T_e(s,1)=- T_{\wbar e}(s,0)
\end{equation}
and $m_E(e)=m_E(\overline{e})$, we have that
\begin{align*}
\sum_{e\,\in\,E}m_E(e)G\left(V_e, T_e\right)(s,1) &=\sum_{e\,\in\,E}m_E(\overline{e})G\left( V_{\overline{e}}, -T_{\wbar e}\right)(s,0)\\
&=-\sum_{e\,\in\,E}m_E(e)G\left( V_e, T_e\right)(s,0),
\end{align*}
and thus,
\begin{equation}\label{f2}
\sum_{e\,\in\,E}m_E(e)\big[G\left(V_e, T_e\right)\big]_{t=0}^{t=1}=-2\sum_{e\,\in\,E}m_E(e)G\left(V_e, T_e\right)(s,0).
\end{equation}
Combining~\eqref{f2} with~\eqref{f1}, we obtain the first variation formula~\eqref{Eq:firstVF}.


Next, we have that
\begin{multline}\label{f3}
\frac{d^2}{d s^2} E_G(f_{s})\\
\begin{aligned}
&=\frac{1}{2}\sum_{e\,\in\,E}m_E(e)\int_0^1 \frac{\p^2}{\p s^2}G\left(T_e, T_e\right)\, dt\\
&=\sum_{e\,\in\,E}m_E(e)\int_0^1 G\left(\nabla_{V_e}\nabla_{V_e}T_e, T_e\right)+\left\|\nabla_{V_e}T_e\right\|^2\,dt \\
&=\sum_{e\,\in\,E}m_E(e)\int_0^1 G\left(\nabla_{V_e}\nabla_{T_e}V_e, T_e\right)+\left\|\nabla_{T_e}V_e\right\|^2\,dt \quad \left(\text{use}\; [V_e,T_e]=0\right)\\
&=\sum_{e\,\in\,E}m_E(e)\int_0^1 G\left(\nabla_{T_e}\nabla_{V_e}V_e, T_e\right)+G\left(R\left(V_e, T_e\right)V_e, T_e\right)+\left\|\nabla_{T_e}V_e\right\|^2\,dt.
\end{aligned}
\end{multline}
The first term becomes
\begin{multline}\label{f4}
\sum_{e\,\in\,E}m_E(e)\int_0^1 G\left(\nabla_{T_e}\nabla_{V_e}V_e, T_e\right)\, dt\\
\begin{aligned}
&=\sum_{e\,\in\,E}m_E(e)\int_0^1 \frac{\p}{\p t}G\left(\nabla_{V_e}V_e, T_e\right)-G\left(\nabla_{V_e}V_e, \nabla_{T_e}T_e\right)\, dt \\
&=-2\sum_{e\,\in\,E}m_E(e)G\left(\nabla_{V_e}V_e, T_e\right)(s,0)-\sum_{e\,\in\,E}m_E(e)\int_0^1G\left(\nabla_{V_e}V_e, \nabla_{T_e}T_e\right)\, dt,
\end{aligned}
\end{multline}
where we used~\eqref{f2} and $\nabla_{V_e}V_e(s,1)=\nabla_{V_{\wbar e}}V_{\wbar e}(s,0)$. Substituting~\eqref{f4} to~\eqref{f3} and using the fact $G(R(V_e, T_e)V_e, T_e)=-G(R(V_e, T_e)T_e, V_e)$, we obtain the second variation formula~\eqref{Eq:secondVF}.
\end{proof}


By the first variation formula~\eqref{Eq:firstVF}, a map $f$ is a critical point of the energy functional $E_G$ if and only if
\begin{equation}\label{harmmap}
\nabla_{T_e}T_e=0\quad \text{and}\quad \sum_{e\,\in\,E_x} m_E(e)T_e(x)=0
\end{equation}
for any edge $e\in E$ and $x\in V$. We call the map satisfying~\eqref{harmmap} a harmonic map. Moreover, according to the second variation formula~\eqref{Eq:secondVF}, we see the Hessian of $E_G$ for a harmonic map $h: X\to (M,G)$ is given by
\[
{\Hess}_{E_G}(V,W)_h=\sum_{e\,\in\,E} m_E(e)\int_0^1 \big\{G\left(\nabla_{T_e}V_e, \nabla_{T_e}W_e\right) -G\left(R\left(V_e, T_e\right)T_e, W_e\right)\big\}\,dt,
\]
for any two variation vector fields $V, W$ of $h$. From the expression, we see that ${\Hess}_{E_G}$ is non-negative definite if $(M,G)$ has non-positive sectional curvature. If furthermore, $(M,G)$ is a compact Riemannian manifold of negative sectional curvature and the image of the harmonic map $h: X\to M$ is not homotopic to a point nor a closed circle, then ${\Hess}_{E_G}$ is non-degenerate at $h$ (see the proof of Theorem~\ref{Thm:KSharmonic}$\MK$\eqref{theo2.1.3}).



\section{Smooth dependence of harmonic maps}\label{AppB}

Let $(M,G)$ be a compact smooth Riemannian manifold. For any integer $k\geq 1$, we denote the set of all $C^k$-Riemannian metrics on $M$ by $\Met^k(M)$. Note that the space $\Met^k(M)$ is an open subset of the Banach space consisting of $C^k$-symmetric $(0,2)$-tensors on $M$.

\begin{prop}\label{smooth dependence}
Fix any integer $k \ge 1$. Let $X$ be a finite weighted graph, $(M, G)$ be a closed Riemannian manifold of nonpositive sectional curvature and $h: X\to M$ be a harmonic map. If the Hessian $\Hess_{E_G}$ of the energy functional $E_G$ is non-degenerate at $h$, then there exists an open neighborhood $\Uc$ of $G$ in $\Met^{k+1}(G)$ and $C^k$-map $\wt h:\Uc \to C^{k+2}(X, M)$ such that $\wt h(G'):X \to M$ is a harmonic map for each $G' \in \Uc$ and $\wt h(G)=h$.
\end{prop}

There is a corresponding result where the domain and the target are Riemannian manifolds by Eells--Lemaire~\cite[Theorem~3.1]{EL} and Koiso~\cite[Theorem~4.7]{Koiso}. We shall prove by a simple application of the implicit function theorem between Banach spaces as in~\cite{EL}. Since we have not found the argument adapted to our setting, we describe the setup in a self-contained manner in the following.



For any piecewise $C^k$-map $f:X \to M$, let us denote the vector space consisting of $C^k$-vector fields along $f$ by $C^k(X, f^{-1}TM)$. We understand that every $v$ in $C^k(X, f^{-1}TM)$ is in $C^k$ on each interior $(0, 1)$ of edges and continuous on $X$. We define a norm on $C^k(X, f^{-1}TM)$ by
\[
\|v\|_k:=\sum_{0\,\leq\,i\,\leq\,k}\underset{e\,\in\,E_0}\sup\;\sup_{t\,\in\,[0,\, 1]}\left\|\nabla^i_{\p_t f_e} v_e\right\|_{G},
\]
where $E_0$ is the set of unoriented edges, for each edge $e\in E_0$, $v_e$ is the restriction of $v$ and $\nabla^i_{\p_t f_e}$ denotes the $i^{\rm th}$ covariant derivative along $f_e$ relative to the metric $G$. Then, $(C^k(X, f^{-1}TM), \|\cdot\|_k)$ becomes a Banach space.


Let $C^k(X, M):=\{f:X\to M: f\ \text{ is a piecewise $C^k$-map}\}$. Then, $C^k(X, M)$ is a Banach manifold modeled on the Banach space $(C^k(X, f^{-1}TM), \|\cdot \|_k)$. Namely, for any $f\in C^k(X, M)$ and an open neighborhood $B\subset C^k(X, f^{-1}TM)$ of $0$, we define the \emph{exponential map} ${\exp}_f:B\to C^k(X,M)$ by $ (\exp_fv)(x):=\exp_{f(x)}v_{x}, $ where $\exp_{f(x)}: T_{f(x)}M\to M$ denotes the exponential map of the Riemannian manifold $(M,G)$. Then, for an open neighborhood $B$, the map $\exp_f^{-1}: \exp_f(B)\to B$ gives a local chart around the point $f$.



We consider the subset of $C^k(X, M)$ consisting of piecewise $C^k$-maps satisfying the balanced condition:
\[
C^k_{bal}(X,M):=\left\{f\in C^k(X,M): \sum_{e\,\in\,E_x}m_E(e)\p_t f_e (x)=0 \quad \text{for any }x\in V\right\},
\]
where $\p_t f_e=(d/dt)f_e$. Note that the set $C^k_{bal}(X,M)$ is not empty since there always exists a harmonic map, which satisfies the balanced condition. For any $f \in C^k_{bal}(X,M)$, let us define the Banach space
\[
\Vc_{f,\,G}^k:=\!\left\{v\in C^k\left(X, f^{-1}TM\right)\,:\, \nabla^G_{v_x}\cdot \left(\sum_{e\,\in\,E_{x}}m_E(e) \p_t{f}_e(x)\right)=0 \quad \text{for any }\;x\in V\right\},
\]
as a closed subspace of $C^k(X, f^{-1}TM)$. Note that the condition in the definition $\Vc_{f,\,G}^k$ is equivalent to
\[
\sum_{e\,\in\,E_{x}}m_E(e) \nabla^G_{\p_t{f}_e}v_e(x)=0 \quad \text{for any }x\in V
\]
since $[v_e, \p_t{f}_e]=0$ for any $e \in E$.


\begin{lemm}
For any integer $k \ge 1$, $C^k_{bal}(X,M)$ is a submanifold of $C^k(X,M)$ and for any $C^k$-metric $G$ on $M$ and any $f \in C^k_{bal}(X,M)$, the tangent space of $C^k_{bal}(X,M)$ at $f$ is isomorphic to the Banach space $\Vc_{f,\,G}^k$.
\end{lemm}


\begin{proof}
Fix an arbitrary $f\in C^k_{bal}(X, M)$. Since $\Vc_{f,\,G}^k$ is a closed subspace of $C^k(X, f^{-1}TM)$, it becomes a Banach space with the induced norm $\|\cdot\|_k$. For a sufficiently small $v\in C^k(X, f^{-1}TM)$ with respect to $\|\cdot\|_k$, we set $f^v:=\exp_f\, v\in C^k(X, M)$. First, we show $f^v\in C^k_{bal}(X, M)$ if $v\in \Vc_{f,\,G}^k$. We take a geodesic normal coordinate $(U,\{x^1,\,\ldots,\,x^m\})$ of $(M,G)$ around a point $f(x)$ for $x\in V$, and define
\[
\exp: U\times \R^m\to M\text{ by } \exp\left(p, v^1\,\ldots,\,v^m\right)=\exp_{p}\left(\sum_{i=1}^mv^i\frac{\p}{\p x^i}\middle|_{p}\right).
\]
By definition of $\exp$ and $f^v$, we have
\[
\frac{d}{d t} {f}^v_e(x)=\frac{d}{dt} \exp\left(f_e(t), v_e(t)\right)\Big|_{t=0}=\left(d\exp\right)_{\left(f(x), v_x\right)}\left(\p_t f_e(x), \nabla^G_{\p_t{f}_e}v_e(x)\right).
\]
for each $e\in E_x$, and thus we obtain
\begin{equation}\label{dexp}
\sum_{e\,\in\,E_x}m_E(e)\frac{d}{d t} f^v_e(x)=\left(d{\exp}\right)_{\left(f(x), \,v_x\right)}\left(0, \sum_{e\,\in\,E_x}m_E(e)\nabla^G_{\p_t{f}_e}v_e(x)\right)
\end{equation}
since $f\in C^k_{bal}(X, M)$. Therefore, if $v\in \Vc_{f,\,G}^k$, then $f^v\in C^k_{bal}(X, M)$ as required.

Thus, we obtain a well-defined map
\[
\exp_f: \Uc\cap \Vc_{f,\,G}^k\to \exp_f(\Uc)\cap C^k_{bal}(X,M)
\]
for an open neighborhood $\Uc$ in $C^k(X, f^{-1}TM)$ around $0$, and this map is injective if $\Uc$ is small enough since it is a restriction of the exponential map $\exp_f$ to $\Vc_{f,\,G}^k$. We shall show the map is surjective, i.e., for any $f^v\in \exp_f(\Uc)\cap C^k_{bal}(X, M)$, the corresponding vector field $v:=\exp^{-1}_f(f^v)\in C^k(X, f^{-1}TM)$ is actually contained in $\Vc_{f, G}^k$. By~\eqref{dexp}, it is sufficient to show that $(d\exp)_{(x,\,v_x)}(0,w)=0$ implies $w=0$. Indeed, we have
\[
(d\exp)_{\left(f(x),\,v_x\right)}(0,w)=d\left(\exp_{f(x)}\right)_{v_x}(w),
\]
and hence, if $\Uc$ is small enough so that for any $v \in \Uc$ satisfies that $\|v\|_k<\epsilon$ for sufficiently small $\epsilon>0$ (for instance, we may take $\epsilon$ as ${\inj}(M)>0$), then $d({\exp}_x)_{v_x}$ is injective, and $w=0$; as required. This implies the lemma.
\end{proof}


We define the map
\begin{gather*}
\t: C_{bal}^{k+2}(X, M) \times \Met^{k+1}(M) \to \prod_{e\,\in\,E_0}\left(C^k\left([0, 1], f_e^{-1}TM\right)\right),
\\
\t(f, G):=\left(\nabla_{\p_t f_e}^G \p_t f_e\right)_{e\,\in\,E_0},
\end{gather*}
where $E_0$ is the set of unoriented edges of $X$; for each $e \in E_0$, we identify $e$ with $[0, 1]$ and denote by $C^k([0, 1], f_e^{-1}TM)$ the space of $C^k$-vector fields along $f_e:[0, 1] \to M$.

Let us fix a piecewise smooth harmonic map $h: X\to (M,G)$. Taking a small enough open neighborhood $U_{bal}\subset \Vc_{h,\,G}^{k+2}$ of $0$, and we identify ${\exp}_h(U_{bal})\subset C^{k+2}_{bal}\linebreak(X, M)$ and $U_{bal}$. For any $f \in U_{bal}$, we identify
\[
\prod_{e\,\in\,E_0}\left(C^k\left([0, 1], f_e^{-1}TM\right)\right) \quad \text{and} \quad \prod_{e\,\in\, E_0}\left(C^k\left([0, 1], h_e^{-1}TM\right)\right)
\]
via parallel transport relative to $\nabla^G$ along the curve $\gamma_f(t,x):=\exp_{h(x)}(tv(x))$ where $f=\exp_h\, v\in U_{bal}$.



We take any smooth curve $s \mapsto (f_s, G_s)$ in $U_{bal}\times \Met^{k+2}(M)$ through $(h, G)$ at $s=0$ for $s \in (-\e, \e)$ and for $\e>0$. Let us write $\t(f_s, G_s)=(\t_{s,\,e})_{e\,\in\,E_0}$ for $s \in (-\e, \e)$. Then, by the equations~\eqref{c4} and~\eqref{c41} given in the proof of Lemma~\ref{le3} (Note that the equations~\eqref{c4} and~\eqref{c41} hold for any Riemannian manifold as a target manifold), we see on each $e \in E_0$,
\begin{equation}\label{derivtau}
\frac{d\tau_{s,\,e}}{ds}\left|_{s=0}=\nabla_{T_e}\nabla_{T_e}V_e+R\left(V_e, T_e\right)T_e+\frac{d}{ds}\right|_{s=0} \nabla^{G_s}_{T_e}T_e,
\end{equation}
where $T_e:=\p_t h_e$ and $V_e:=(\partial_s f_{s, \,e})|_{s=0}$. Note that the final term in the right hand side depends only on $G_s$ and the initial condition $h$.

\begin{proof}[Proof of Proposition~\ref{smooth dependence}]
For a harmonic map $h:X \to (M, G)$, we take an open neighborhood $U_{bal}$ of $h$ in $\Vc_{h,\,G}^{k+2}$ and an open neighborhood $\Uc\subset \Met^{k+1}(M)$ of $G$. We regard $\tau$ as a $C^k$-map between Banach spaces:
\[
\tau: U_{bal}\times \Uc\to \prod_{e\,\in\,E_0}\left(C^k\left([0, 1], h_e^{-1}TM\right)\right).
\]
We shall show that the derivative of $\t$ in the $U_{bal}$-component at $o=(h, G)$
\[
d\tau|_{U_{bal},\,o}: \Vc_{h,\,G}^{k+2}\to \prod_{e\,\in\,E_0}\left(C^k\left([0, 1], h_e^{-1}TM\right)\right)
\]
is an isomorphism.

By~\eqref{derivtau}, we have that $d\tau_e|_{U_{bal},\,o}(V)=\nabla_{T_e}\nabla_{T_e}V_e+R(V_e, T_e)T_e$ on each $e \in E_0$ for any $V\in \Vc_{h,\, G}^{k+2}$ since the last term in~\eqref{derivtau} is independent of the deformation $f_s$ of $h$ relative to $V$. Thus, $d\tau|_{U_{bal},\,o}$ is a bounded linear map. First we show that $d\tau|_{U_{bal},\,o}$ is injective. Indeed, if $d\tau|_{U_{bal},\,o}(V)=0$, then taking the deformation $f_s$ relative to $V$ in $C^{k+2}_{bal}(X, M)$ for any $s\in (-\e, \e)$ for $\e>0$, we see that
\begin{align*}
0&=\sum_{e\,\in\,E}m_E(e)\int_0^1 G\left(-\nabla_{T_e}\nabla_{T_e}V_e-R\left(V_e, T_e\right)T_e, V_e\right)\, dt\\
&\quad+\sum_{e\,\in\,E}m_E(e)G_0\left(\nabla_{T_e}V_e, V_e\right)\Big|_{t=0}^{t=1}\\
&={\Hess}_{E_G}(V,V)_h,
\end{align*}
where we used Lemma~\ref{le2} in the first equality which follows from the assumption $f_s\in C^{k+2}_{bal}(X,M)$. The assumption in the statement implies that ${\Hess}_{E_G}$ is non-degenerate at $h$, we obtain $V=0$, and thus $d\tau|_{U_{bal},\,o}$ is injective.

Next we show that $d\tau|_{U_{bal},\,o}$ is surjective. For any given $W=(W_e)_{e \in E_0}$ where $W_e \in C^k([0, 1], h_e^{-1}TM)$ for each $e \in E_0$, we consider the second order ordinary differential equations
\begin{equation}\label{dtauw}
\nabla_{T_e}\nabla_{T_e}V_e+R(V_e, T_e)T_e=W_e
\end{equation}
on each $e \in E_0$ with the condition that $V$ is in $\Vc_{h,\,G}^{k+2}$, i.e., $V$ solves~\eqref{dtauw} on each interior $(0, 1)$ of edges $e$, and is continuous on $X$ satisfying the balanced condition. Let us consider the Hilbert space $H^{1, \,2}_{h,\,G}$ as the completion of $\Vc_{h,\,G}^{k+2}$ endowed with the inner product
\[
\left\langle V^1, V^2\right\rangle_{H_{h,\,G}^{1,\,2}}:=\sum_{e \,\in\,E}m_E(e)\int_0^1\Big\{G\left(\nabla_{T_e}V_e^1, \nabla_{T_e}V_e^2\right)-G\left(R\left(V_e^1, T_e\right)T_e, V_e^2\right)\Big\}\,dt,
\]
where this indeed defines the positive definite inner product since $(M, G)$ has nonpositive sectional curvature and ${\Hess}_{E_G}$ is non-degenerate at $h$. Note that every $V$ in $H^{1,\,2}_{h,\,G}$ is continuous on $X$. For any $W \in \prod_{e\,\in\,E_0}(C^k([0, 1], h_e^{-1}TM))$, we define a linear functional $L_W$ on $H^{1,\,2}_{h,\, G}$ by

\[
L_W(\f):=\sum_{e\,\in\,E}m_E(e)\int_0^1 G\left(W_e, \f_e\right)\,dt, \quad \text{for }\f \in H^{1,\,2}_{h,\,G}.
\]
Then, $L_W$ is a bounded linear functional on $H^{1, 2}_{h,\,G}$; in fact, the natural embedding map $H^{1,\,2}_{h,\,G} \to C^0(X, h^{-1}TM)$ is compact. Therefore the Riesz representation theorem implies that there exists a unique $V$ in $H^{1,\,2}_{h,\,G}$ such that $\langle V, \f\rangle_{H_{h,\,G}^{1,\,2}}=-L_W(\f)$ for any $\f$ in $H^{1,\,2}_{h,\,G}$. This implies that $V$ solves~\eqref{dtauw} on each interior $(0, 1)$ of edges $e$ and $V_e$ is in $C^{k+2}$ on $(0, 1)$ by taking smooth functions $\f_e$ whose supports are included in $(0, 1)$ for each $e \in E_0$. Then, integration by parts gives
\begin{multline*}
\sum_{e\,\in\,E}m_E(e)\int_0^1\big\{G\left(-\nabla_{T_e}\nabla_{T_e}V_e-R\left(V_e, T_e\right)T_e, \f_e\right)\big\}\,dt -2\sum_{x\,\in\,V}\sum_{e\,\in\,E_x}G\left(\nabla_{T_e}V_e, \f_e\right)\\
=-\sum_{e\,\in\,E}m_E(e)\int_0^1 G\left(W_e, \f_e\right)\,dt,
\end{multline*}
and $\sum_{x\,\in\,V}\sum_{e\,\in\,E_x}G(\nabla_{T_e}V_e, \f_e)=0$ holds for any $\f$ in $H^{1,\,2}_{h,\,G}$. Therefore $V$ satisfies the balanced condition and is in $V_{h,\,G}^{k+2}$. Hence $d\tau|_{U_{bal},\,o}$ is surjective.



Now, the implicit function theorem for Banach spaces implies that there exist an open neighborhood $\Uc'\subset \Uc$ of $G$ and a unique $C^k$-map $\wt h: \Uc'\to U_{bal}$ satisfying that
\[
\tau\left(\wt h\left(G'\right),G'\right)=0
\]
on $\Uc'$. Then, $\wt h(G): X\to M$ is a harmonic map relative to the metric $G'$ for any $G' \in \Uc'$ and $\wt h(G)=h$, and we complete the proof.
\end{proof}

